\chapter{Nonholonomic Constraints}
\label{ch:nonholonomic}

A constraint of the form $f(q,t)=0$, where $q=(q^1,\dots,q^n)$ are
generalized coordinates, is called \emph{holonomic}: it can be
solved, at least locally, to eliminate one coordinate outright, and
it reduces the dimension of the configuration space itself. A more
general class of constraints is linear in the velocities,
\begin{equation}
  \sum_{i=1}^n a_i(q,t) \dot q^i + a_t(q,t) = 0,
  \label{pfaffian}
\end{equation}
a so-called Pfaffian constraint. If the left-hand side of
\eqref{pfaffian} happens to be the total time derivative of some
function $f(q,t)$, the constraint is secretly holonomic in disguise.
If it is not --- if no such $f$ exists --- the constraint is called
\emph{nonholonomic}: it restricts the \emph{velocities} available at
each point of $Q$, without restricting which \emph{configurations}
are reachable.

The canonical example is a coin (or disk) of radius $R$ rolling
without slipping on a plane, with contact point $(x,y)$, heading
angle $\phi$, and roll angle $\theta$ as coordinates. The
rolling condition reads
\begin{equation*}
  \dot x = R\,\dot\theta\cos\phi ,\qquad
  \dot y = R\,\dot\theta\sin\phi ,
\end{equation*}
which is not integrable: it constrains the instantaneous velocity to
a two-dimensional subspace of the four-dimensional tangent space at
each point, and yet, by alternating steering and rolling, the coin
can be driven from \emph{any} configuration $(x,y,\phi)$ to any
other. The instantaneous constraint is genuine, but it places no
restriction whatsoever on the reachable set of configurations --- an
instance of what control theorists call bracket-generating, or
Chow--Rashevski\u{i}, controllability. This already runs against a
piece of intuition carried over uncritically from the holonomic
case, where a nontrivial velocity constraint always shrinks the
accessible configuration space too.

A second standard example, the \emph{Chaplygin sleigh}, makes an
even sharper point.

It is a rigid body sliding frictionlessly on a
plane, with a knife edge that permits motion only along its own
length (no sideways slip). The system conserves energy exactly ---
there is no friction anywhere in the model --- and yet its reduced
dynamics on the velocity variables is attracted to a fixed
equilibrium: left alone, the sleigh asymptotically settles into
straight-line motion with no rotation, knife edge trailing the
center of mass~\cite{chaplygin}. This
happens because the reduced equations of motion on the nonholonomic
constraint surface, unlike those of an ordinary holonomic mechanical
system, do not in general preserve phase-space (Liouville) volume,
even though they do preserve energy. A perfectly conservative
mechanical system can therefore display genuinely dissipative-looking
long-time behaviour --- attracting fixed points, basins of
attraction --- with no friction anywhere in the equations. Energy
conservation, by itself, is not enough to rule out the qualitative
signatures usually associated with dissipation once constraints stop
being holonomic.

For holonomic constraints, two routes to the equations of motion ---
D'Alembert's principle (constraint forces do no virtual work, for
virtual displacements compatible with the instantaneous constraint)
and Hamilton's principle (the true path extremizes the action among
\emph{all} paths satisfying the constraint) --- are provably
equivalent. For nonholonomic (non-integrable, velocity-linear)
constraints, this equivalence simply fails, and the failure is not a
minor technicality.

There are two natural-looking ways to generalize the variational
treatment to a system with a nonholonomic constraint of the
form~\eqref{pfaffian}.

\begin{description}
  \item[Nonholonomic (Lagrange--d'Alembert, or Chetaev) mechanics.]
  Impose the constraint directly on virtual displacements at each
  instant, exactly as in the holonomic case, and only afterwards
  derive equations of motion:
  \begin{equation}
    \frac{d}{dt}\frac{\partial L}{\partial \dot q^i}
    - \frac{\partial L}{\partial q^i}
    \;=\; \sum_a \lambda_a\, a_i^a(q),
    \qquad
    \sum_i a_i^a(q)\,\dot q^i + a_t^a(q) = 0 .
    \label{lagrange-dalembert}
  \end{equation}
  \item[Vakonomic mechanics.] (The name --- "variational axiomatic
  kind" --- was coined by Kozlov~\cite{kozlov}.)
  Instead, treat the nonholonomic constraint as a genuine constraint
  on the space of \emph{paths} from the outset, as one would an
  isoperimetric constraint in the calculus of variations, and only
  then extremize the action over that constrained set of paths.
\end{description}

For holonomic constraints these two prescriptions coincide. For
genuinely nonholonomic constraints they generically do not: it has
been shown, for instance, that for the rolling coin on an inclined
plane no nontrivial vakonomic solution reproduces the
Lagrange--d'Alembert solution at
all~\cite{zampieri}. This is not merely a
mathematical curiosity to be settled by taste: Lewis and Murray
built and instrumented a physical system --- a ball rolling without
slipping on a rotating table --- specifically to adjudicate between
the two theories, and found that the observed motion matches the
Lagrange--d'Alembert (nonholonomic) prediction; they were unable to
produce vakonomic simulations resembling the experimental data at
all~\cite{lewis+murray}. The variational principle that looks like
the most direct generalization of Hamilton's principle to velocity
constraints is, for constraints realized by an actual rolling or
sliding contact, simply the wrong physical theory.

This does not make vakonomic mechanics useless --- it correctly
describes optimal-control and sub-Riemannian-geodesic problems, where
the "constraint" genuinely is a restriction on the class of paths
being optimized over, rather than a constraint enforced
instant-by-instant by a physical reaction force. The moral is
subtler and more interesting than "one theory is right and the other
wrong": the bare constraint equation~\eqref{pfaffian}, by itself,
does not determine the dynamics. The correct equations of motion
depend on how the constraint is physically \emph{realized} ---
information that is invisible at the level of the constraint
equation alone, and that a first course in mechanics, built entirely
on holonomic examples where this ambiguity cannot arise, gives no
reason to expect.


\paragraph{Nonholonomic constraints: disagreement}
\label{nonholonomic}

A nonholonomic constraint restricts velocities, not positions:
\begin{equation}
  a^k_i(q)\,\dot q^i=0,\qquad k=1,\dots,m,
  \label{nonhol}
\end{equation}
where the one-forms $a^k=a^k_i\,dq^i$ do not integrate to the form $df_k$.
Two natural generalizations of the holonomic procedure exist, and they give
different equations (see subsection \ref{ap:nh}) .

\subparagraph{Route A: Lagrange--d'Alembert (Newton-first).}
Take~\eqref{ideal} as the definition of an ideal constraint and impose the
constraint on virtual displacements:
$a^k_i\,\delta q^i=0$. Then
$\bigl(\frac{d}{dt}\frac{\partial L}{\partial\dot q^i}-\frac{\partial L}{\partial q^i}\bigr)\delta q^i=0$
for all admissible $\delta q$, so
\begin{equation}
  \frac{d}{dt}\frac{\partial L}{\partial\dot q^i}-\frac{\partial L}{\partial q^i}
  =\lambda_k\,a^k_i .
  \label{LdA}
\end{equation}
This is \emph{not} the Euler--Lagrange equation of any augmented action built
from~\eqref{nonhol}; it is a virtual-work principle that does not
integrate into~\eqref{hamilton}.

\subparagraph{Route B: vakonomic (action-first).}
Treat the problem as a variational problem with a nonholonomic side
condition. Extremize $S$ among paths satisfying~\eqref{nonhol}; the
multipliers enter as in~\eqref{aug-holonomic},
$\tilde L=L+\lambda_k a^k_i\dot q^i$, and the Euler--Lagrange equations
for $q^i$ are
\begin{equation}
  \frac{d}{dt}\frac{\partial L}{\partial\dot q^i}-\frac{\partial L}{\partial q^i}
  =-\dot\lambda_k\,a^k_i-\lambda_k\bigl(\partial_ja^k_i-\partial_ia^k_j\bigr)\dot q^j .
  \label{vako}
\end{equation}

Comparing~\eqref{LdA} and~\eqref{vako}: the term $-\dot\lambda_ka^k_i$
alone could be absorbed by redefining the multiplier, which is what happens
in the holonomic case. The second term involves the \emph{curvature} of the
constraint distribution, $d a^k=\tfrac12(\partial_ja^k_i-\partial_ia^k_j)\,dq^j\wedge dq^i$,
and is nonzero precisely when the constraint is not integrable. In other
words:
\begin{quote}
\emph{The two procedures coincide if and only if the constraint is
holonomic.}
\end{quote}
The disagreement is not a computational subtlety. It is a statement about
how nature behaves, and nature has voted against vakonomic (subsection \ref{ap:nh}).

\begin{example}[A skate]\label{skate}
A skate is a rigid body on a plane whose blade touches the ice at its center
of mass. Coordinates $(x,y,\varphi)$, with Lagrangian
$L=\tfrac12m(\dot x^2+\dot y^2)+\tfrac12I\dot\varphi^2$, and the blade cannot
slip sideways:
\[
  a=\sin\varphi\,dx-\cos\varphi\,dy,\qquad a_\varphi=0 .
\]
\emph{Route~A.} The reaction is perpendicular to the blade and applied at the
center of mass, so it exerts no torque: $I\ddot\varphi=0$ and the speed along
the blade is constant. This is what a skater coasting on a straight line
observes.

\emph{Route~B.} The $\varphi$-component of~\eqref{vako} reads
\[
  I\ddot\varphi=\lambda\,v,\qquad v=\dot x\cos\varphi+\dot y\sin\varphi,
\]
so in general the ``constraint'' exerts a torque about the center of mass
which the real blade cannot exert. Moreover, Route~B is a
\emph{boundary-value} problem: for given endpoints the multiplier is fixed
by the final configuration, and the particle's initial velocity can no longer be
prescribed freely.
\end{example}

Experiments on rolling bodies, skates and sleighs (and the engineering
practice of wheeled vehicles) are consistent with Route~A. Route~B is not
wrong, but it answers a different question: it describes the optimal control
of a system whose velocities are \emph{steered} (a sub-Riemannian
geodesic problem, with applications to robotics and parking a car), not the
free motion of a system whose velocities are \emph{forced} by a passive
constraint. The principle of least action is therefore a theorem about
\emph{potential} dynamics under \emph{holonomic} constraints, and a
definition or a hypothesis beyond that.

